%%%%%%%%%%%%%%%%
%%% deut 24 %%%
%%%%%%%%%%%%%%%%

\Chapter{24}

\Verse{1}
{
	\hDtn{כִּי יִקַּח אִישׁ אִשָּׁה וּבְעָלָהּ וְהָיָה אִם לֹא תִמְצָא חֵן בְּעֵינָיו כִּי מָצָא בָהּ עֶרְוַת דָּבָר וְכָתַב לָהּ סֵפֶר כְּרִיתֻת וְנָתַן בְּיָדָהּ וְשִׁלְּחָהּ מִבֵּיתוֹ׃}
}
{
	\eDtn{
		24 “When a man will take a woman and marry her, and it\\
		will be that, if she does not find favor in his eyes\\
		because he has found an exposure of something in her, and\\
		he will write a document of cutting-off for her and put it\\
		in her hand and let her go from his house,\\
	}
}
{
%	\footnote{
%		When a man will take a woman. This law (vv. 1–4) has been taken as the
%		biblical law of divorce, but it is not. It is the law governing a
%		specific instance in which a couple might want to return to each other
%		after they were divorced and she was remarried and then was divorced
%		again or widowed. Divorce law in general has been derived in part from
%		this case because of the curious fact that there is no law in the Torah
%		telling how to get married and no law telling how to get divorced. Here,
%		the divorce procedure appears to be assumed, as a known practice. The
%		absence of marriage and divorce procedures has been used as a proof of
%		the existence of oral law that was given along with the law that is
%		written in the Torah. But the case may be rather that these ceremonies
%		were not regarded as having been given in detail by God—as few
%		ceremonies outside of the sacrifices at the Tabernacle are so given. The
%		divine interest is in the marriage relationship rather than in the
%		ceremony: the husband’s treatment of the wife, the prohibition of
%		adultery, their relationship with their children, inheritance, making
%		one’s spouse happy (see 24:5). The same goes for funerals and bar
%		mitzvah. Even such well-known Jewish practices as saying kaddish and
%		sitting shivah are not mentioned in the Torah. And bar and bat mitzvah
%		are not treated (or even mentioned): the concern is with doing the
%		commandments. The details of the ceremony of arriving at the age of
%		doing them are left for humans to choose. All of these ceremonies were
%		detailed in rabbinic law and practice later. Ceremonial procedures are
%		given in the Torah for less common cases, such as the slave who declines
%		to leave a master (Deut 15:16–17), or the brother-in-law who declines to
%		perform levirate marriage (25:7–10), or the present case of return
%		marriage. Footnote 24:1. an exposure of something. See the comment on
%		23:15 above. Footnote 24:1. a document of cutting-off. A divorce.
%	}
}

\Verse{2}
{
	\hDtn{וְיָצְאָה מִבֵּיתוֹ וְהָלְכָה וְהָיְתָה לְאִישׁ אַחֵר׃}
}
{
	\eDtn{
		and she will go out from his house and go and become\\
		another man’s,\\
	}
}
{
}

\Verse{3}
{
	\hDtn{וּשְׂנֵאָהּ הָאִישׁ הָאַחֲרוֹן וְכָתַב לָהּ סֵפֶר כְּרִיתֻת וְנָתַן בְּיָדָהּ וְשִׁלְּחָהּ מִבֵּיתוֹ אוֹ כִי יָמוּת הָאִישׁ הָאַחֲרוֹן אֲשֶׁר לְקָחָהּ לוֹ לְאִשָּׁה׃}
}
{
	\eDtn{
		and the latter man will hate her and write a document of\\
		cutting-off for her and put it in her hand and let her go\\
		from his house, or if the latter man who took her to him\\
		for a wife will die:\\
	}
}
{
}

\Verse{4}
{
	\hDtn{לֹא יוּכַל בַּעְלָהּ הָרִאשׁוֹן אֲשֶׁר שִׁלְּחָהּ לָשׁוּב לְקַחְתָּהּ לִהְיוֹת לוֹ לְאִשָּׁה אַחֲרֵי אֲשֶׁר הֻטַּמָּאָה כִּי תוֹעֵבָה הִוא לִפְנֵי יְהֹוָה וְלֹא תַחֲטִיא אֶת הָאָרֶץ אֲשֶׁר יְהֹוָה אֱלֹהֶיךָ נֹתֵן לְךָ נַחֲלָה׃}
}
{
	\eDtn{
		her first husband who let her go shall not be able to come\\
		back to take her to be his for a wife since she has been\\
		made impure, because that is an offensive thing in front\\
		of YHWH, and you shall not bring sin on the land that\\
		YHWH, your God, is giving you as a legacy.\\
	}
}
{
%	\footnote{
%		her first husband shall not be able to come back. Jeremiah transforms
%		this into a metaphor in which God has cut off Israel the way the husband
%		in this law cuts off his wife (Jer 3:1), yet God would still take Israel
%		back, seemingly against God’s own law concerning human marriages
%		(3:12,14,22)!
%	}
}

\Verse{5}
{
	\hDtn{כִּי יִקַּח אִישׁ אִשָּׁה חֲדָשָׁה לֹא יֵצֵא בַּצָּבָא וְלֹא יַעֲבֹר עָלָיו לְכל דָּבָר נָקִי יִהְיֶה לְבֵיתוֹ שָׁנָה אֶחָת וְשִׂמַּח אֶת אִשְׁתּוֹ אֲשֶׁר לָקָח׃}
}
{
	\eDtn{
		“When a man will take a new wife, he shall not go out in\\
		the army and not go along with it for any matter. He shall\\
		be free at his house for one year and shall make his wife\\
		whom he has taken happy.\\
	}
}
{
}

\Verse{6}
{
	\hDtn{לֹא יַחֲבֹל רֵחַיִם וָרָכֶב כִּי נֶפֶשׁ הוּא חֹבֵל׃}
}
{
	\eDtn{
		One shall not take a mill or an upper millstone as\\
		security, because he is taking one’s life as security.\\
	}
}
{
}

\Verse{7}
{
	\hDtn{כִּי יִמָּצֵא אִישׁ גֹּנֵב נֶפֶשׁ מֵאֶחָיו מִבְּנֵי יִשְׂרָאֵל וְהִתְעַמֶּר בּוֹ וּמְכָרוֹ וּמֵת הַגַּנָּב הַהוּא וּבִעַרְתָּ הָרָע מִקִּרְבֶּךָ׃}
}
{
	\eDtn{
		“When a man will be found stealing a person from among his\\
		brothers, from the children of Israel, so he will get\\
		profit through him and sell him, then that thief shall\\
		die. So you shall burn away what is bad from among you.\\
	}
}
{
}

\Verse{8}
{
	\hDtn{הִשָּׁמֶר בְּנֶגַע הַצָּרַעַת לִשְׁמֹר מְאֹד וְלַעֲשׂוֹת כְּכֹל אֲשֶׁר יוֹרוּ אֶתְכֶם הַכֹּהֲנִים הַלְוִיִּם כַּאֲשֶׁר צִוִּיתִם תִּשְׁמְרוּ לַעֲשׂוֹת׃}
}
{
	\eDtn{
		“Be watchful with the plague of leprosy, to be very\\
		watchful and to do according to everything that the Levite\\
		priests will instruct you. You shall be watchful to do\\
		according to what I commanded them.\\
	}
}
{
}

\Verse{9}
{
	\hDtn{זָכוֹר אֵת אֲשֶׁר עָשָׂה יְהֹוָה אֱלֹהֶיךָ לְמִרְיָם בַּדֶּרֶךְ בְּצֵאתְכֶם מִמִּצְרָיִם׃}
}
{
	\eDtn{
		Remember what YHWH, your God, did to Miriam on the way\\
		when you were coming out from Egypt.\\
	}
}
{
%	\footnote{
%		Remember what God did to Miriam. Why is this brought up here in the
%		context of a general law about leprosy? Because the law here is not
%		concerned only with the correct treatment for leprosy after one has
%		contracted it (which is the concern of Leviticus 13–14). It also warns
%		to be watchful (the term occurs three times) against behavior that would
%		bring about leprosy as a consequence (see 24:8). Miriam does something
%		that causes her to be stricken with leprosy. Rashi understands her
%		offense to be slanderous gossip (l[image][image]ôn h[image]r[image]‘),
%		but leprosy is more likely to be tied to a ritual offense. Miriam’s
%		words are not gossip; she is stricken with leprosy for having challenged
%		Moses’ status as a prophet compared with her own (Num 12:9–15; see the
%		comment on 12:2). And later, King Uzziah is stricken with it for burning
%		incense in the Temple, which only priests can do (2 Chr 26:16–21).
%	}
}

\Verse{10}
{
	\hDtn{כִּי תַשֶּׁה בְרֵעֲךָ מַשַּׁאת מְאוּמָה לֹא תָבֹא אֶל בֵּיתוֹ לַעֲבֹט עֲבֹטוֹ׃}
}
{
	\eDtn{
		“When you’ll make a loan of anything to your neighbor, you\\
		shall not come into his house to get his pledge.\\
	}
}
{
}

\Verse{11}
{
	\hDtn{בַּחוּץ תַּעֲמֹד וְהָאִישׁ אֲשֶׁר אַתָּה נֹשֶׁה בוֹ יוֹצִיא אֵלֶיךָ אֶת הַעֲבוֹט הַחוּצָה׃}
}
{
	\eDtn{
		You shall stand outside, and the man to whom you’re\\
		lending shall bring the pledge outside to you.\\
	}
}
{
}

\Verse{12}
{
	\hDtn{וְאִם אִישׁ עָנִי הוּא לֹא תִשְׁכַּב בַּעֲבֹטוֹ׃}
}
{
	\eDtn{
		And if he is a poor man, you shall not lie down with his\\
		pledge.\\
	}
}
{
}

\Verse{13}
{
	\hDtn{הָשֵׁב תָּשִׁיב לוֹ אֶת הַעֲבוֹט כְּבוֹא הַשֶּׁמֶשׁ וְשָׁכַב בְּשַׂלְמָתוֹ וּבֵרְכֶךָּ וּלְךָ תִּהְיֶה צְדָקָה לִפְנֵי יְהֹוָה אֱלֹהֶיךָ׃}
}
{
	\eDtn{
		You shall give back the pledge to him as the sun sets, and\\
		he’ll lie down with his clothing, and he’ll bless you, and\\
		you’ll have virtue in front of YHWH, your God.\\
	}
}
{
%	\footnote{
%		you’ll have virtue. See the comment on Deut 6:25.
%	}
}

\Verse{14}
{
	\hDtn{לֹא תַעֲשֹׁק שָׂכִיר עָנִי וְאֶבְיוֹן מֵאַחֶיךָ אוֹ מִגֵּרְךָ אֲשֶׁר בְּאַרְצְךָ בִּשְׁעָרֶיךָ׃}
}
{
	\eDtn{
		“You shall not exploit a poor or an indigent employee,\\
		from your brothers or from your aliens who are in your\\
		land, in your gates.\\
	}
}
{
}

\Verse{15}
{
	\hDtn{בְּיוֹמוֹ תִתֵּן שְׂכָרוֹ וְלֹא תָבוֹא עָלָיו הַשֶּׁמֶשׁ כִּי עָנִי הוּא וְאֵלָיו הוּא נֹשֵׂא אֶת נַפְשׁוֹ וְלֹא יִקְרָא עָלֶיךָ אֶל יְהֹוָה וְהָיָה בְךָ חֵטְא׃}
}
{
	\eDtn{
		You shall give his pay in his day, and the sun shall not\\
		set on it—because he is poor, and he maintains his life by\\
		it—so he won’t call against you to YHWH, and it will be a\\
		sin in you.\\
	}
}
{
}

\Verse{16}
{
	\hDtn{לֹא יוּמְתוּ אָבוֹת עַל בָּנִים וּבָנִים לֹא יוּמְתוּ עַל אָבוֹת אִישׁ בְּחֶטְאוֹ יוּמָתוּ׃}
}
{
	\eDtn{
		“Fathers shall not be put to death for sons, and sons\\
		shall not be put to death for fathers. They shall each be\\
		put to death through his own sin.\\
	}
}
{
%	\footnote{
%		each through his own sin. This does not contradict the statement that
%		God “reckons fathers’ crime on children and on children’s children, on
%		third generations and on fourth generations” (Exod 34:7). That applies
%		to divine justice and may refer to the way in which behavior recurs
%		through generations in a family. This applies to human justice and
%		refers to the point of law that Israelite courts cannot execute people
%		for their relatives’ offenses.
%	}
}

\Verse{17}
{
	\hDtn{לֹא תַטֶּה מִשְׁפַּט גֵּר יָתוֹם וְלֹא תַחֲבֹל בֶּגֶד אַלְמָנָה׃}
}
{
	\eDtn{
		“You shall not bend judgment of an alien or an orphan, and\\
		you shall not take a widow’s clothing as security.\\
	}
}
{
}

\Verse{18}
{
	\hDtn{וְזָכַרְתָּ כִּי עֶבֶד הָיִיתָ בְּמִצְרַיִם וַיִּפְדְּךָ יְהֹוָה אֱלֹהֶיךָ מִשָּׁם עַל כֵּן אָנֹכִי מְצַוְּךָ לַעֲשׂוֹת אֶת הַדָּבָר הַזֶּה׃}
}
{
	\eDtn{
		And you shall remember that you were a slave in Egypt, and\\
		YHWH, your God, redeemed you from there. On account of\\
		this I command you to do this thing.\\
	}
}
{
}

\Verse{19}
{
	\hDtn{כִּי תִקְצֹר קְצִירְךָ בְשָׂדֶךָ וְשָׁכַחְתָּ עֹמֶר בַּשָּׂדֶה לֹא תָשׁוּב לְקַחְתּוֹ לַגֵּר לַיָּתוֹם וְלָאַלְמָנָה יִהְיֶה לְמַעַן יְבָרֶכְךָ יְהֹוָה אֱלֹהֶיךָ בְּכֹל מַעֲשֵׂה יָדֶיךָ׃}
}
{
	\eDtn{
		When you’ll reap your harvest in your field, and you’ll\\
		forget a sheaf in the field, you shall not go back to take\\
		it. It shall be the alien’s and the orphan’s and the\\
		widow’s, so that YHWH, your God, will bless you in all\\
		your hands’ work.\\
	}
}
{
}

\Verse{20}
{
	\hDtn{כִּי תַחְבֹּט זֵיתְךָ לֹא תְפַאֵר אַחֲרֶיךָ לַגֵּר לַיָּתוֹם וְלָאַלְמָנָה יִהְיֶה׃}
}
{
	\eDtn{
		When you’ll beat your olive trees, you shall not do a\\
		bough afterward. It shall be the alien’s and the orphan’s\\
		and the widow’s.\\
	}
}
{
}

\Verse{21}
{
	\hDtn{כִּי תִבְצֹר כַּרְמְךָ לֹא תְעוֹלֵל אַחֲרֶיךָ לַגֵּר לַיָּתוֹם וְלָאַלְמָנָה יִהְיֶה׃}
}
{
	\eDtn{
		When you’ll cut off grapes of your vineyard, you shall not\\
		glean afterward. It shall be the alien’s and the orphan’s\\
		and the widow’s.\\
	}
}
{
}

\Verse{22}
{
	\hDtn{וְזָכַרְתָּ כִּי עֶבֶד הָיִיתָ בְּאֶרֶץ מִצְרָיִם עַל כֵּן אָנֹכִי מְצַוְּךָ לַעֲשׂוֹת אֶת הַדָּבָר הַזֶּה׃}
}
{
	\eDtn{
		And you shall remember that you were a slave in the land\\
		of Egypt. On account of this I command you to do this\\
		thing.\\
	}
}
{
}
